% 09_funcion_utilidad.tex - Capítulo 9: Función de Utilidad (detalle)
\chapter{Función de Utilidad: Detalle y Análisis}
\label{ch:utilidad}

Este capítulo profundiza en la función de utilidad $U(s)$, su diseño, pesos, limitaciones y
cómo explicarla en sustentación.

\section{Formulación Completa}
\label{sec:formulacion-completa}

La utilidad siempre se mide desde la perspectiva de \textbf{MAX = bando que mueve en la raíz}.

\begin{equation}
U(s) = w_{ent} \cdot D_{ent}(s) + w_{car} \cdot C_{car}(s) - w_{dist} \cdot D_{goal}(s) + w_{riv} \cdot D_{riv}(s)
\label{eq:utilidad}
\end{equation}

Donde:
\begin{align}
D_{ent}(s) &= \begin{cases}
+1 & \text{si MAX entregó (ganó)} \\
-1 & \text{si MIN entregó (ganó)} \\
0 & \text{si no hay victoria}
\end{cases} \\
C_{car}(s) &= \begin{cases}
+1 & \text{si MAX porta el recurso} \\
-1 & \text{si MIN porta el recurso} \\
0 & \text{si nadie porta}
\end{cases} \\
D_{goal}(s) &= \text{distancia Manhattan al objetivo actual} \\
D_{riv}(s) &= \begin{cases}
\min_{r \in Rivales} dist(r, pos_{portador}) & \text{si MAX porta} \\
0 & \text{si MIN porta o nadie porta}
\end{cases}
\end{align}

\section{Justificación de Cada Término}
\label{sec:justificacion-terminos}

\begin{table}[H]
\centering
\caption{Justificación de términos de la utilidad}
\begin{tabular}{llp{6cm}}
\toprule
\textbf{Término} & \textbf{Peso} & \textbf{Por qué} \\
\midrule
Entrega ($w_{ent}=10$) & 10 & Victoria/derrota es el objetivo final; debe dominar cualquier combinación de distancias \\
Carga ($w_{car}=3$) & 3 & Portar el recurso es el paso crítico hacia la victoria; vale más que distancia \\
Distancia objetivo ($w_{dist}=1$) & 1 & Cada paso de Manhattan = -1; progreso visible y monótono \\
Distancia rival ($w_{riv}=0.5$) & 0.5 & Rival cerca = riesgo de interceptación; rompe empates entre rutas equidistantes \\
\bottomrule
\end{tabular}
\end{table}
\label{tab:justificacion-utilidad}


\section{Objetivo Dinámico (Distancia al Objetivo)}
\label{sec:objetivo-dinamico}

El objetivo hacia el que se mide la distancia \textbf{cambia según quién porta el recurso}:

\begin{table}[H]
\centering
\caption{Objetivo de la distancia según portador}
\begin{tabular}{lll}
\toprule
\textbf{Situación} & \textbf{Objetivo} & \textbf{Desde dónde se mide} \\
\midrule
MAX porta & Base propia de MAX & Desde el portador \\
MIN porta & Posición del portador (perseguir) & Desde la unidad MAX más cercana \\
Nadie porta & Celda del recurso & Desde la unidad MAX más cercana \\
\bottomrule
\end{tabular}
\end{table}
\label{tab:objetivo-dinamico}


\textbf{Por qué:} Si MIN porta, MAX debe \textit{perseguir al portador} para interceptar o
forzar errores; medir distancia al recurso sería contraproducente (el recurso se mueve con
MIN). Si nadie porta, ambos bandos van al recurso.

\section{Termino Rival: Solo Cuando MAX Porta}
\label{sec:termino-rival}

\begin{equation}
D_{riv}(s) = \begin{cases}
\min_{r \in Unidades_{MIN}} dist(r, pos_{portador_{MAX}}) & \text{si MAX porta} \\
0 & \text{en otro caso}
\end{cases}
\end{equation}

\textbf{Por qué solo cuando MAX porta:} Un rival cerca del portador de MAX = amenaza de
interceptación (pierde el recurso). Si MIN porta, un rival de MAX cerca del portador de MIN
\textit{no} es una amenaza para MAX (es una oportunidad de interceptación, que ya captura la
persecución via $D_{goal}$).

\section{Pesos Configurables en JSON (RNF-03)}
\label{sec:pesos-json}

Los pesos se leen de la sección opcional \code{utilidad} del JSON:

\begin{minted}{json}
{
  "utilidad": {
    "entrega": 10,
    "recurso_cargado": 3,
    "distancia_objetivo": 1,
    "distancia_rival": 0.5
  }
}
\end{minted}

Si falta la clave, se usa el default (10/3/1/0.5). \textbf{Nada está hardcodeado} (RNF-03).

\section{Limite Conocida: Recurso en Celda Propia vs. Llevarlo Lejos}
\label{sec:limitacion-utilidad}

Con pesos 10/3/1/0.5, surge una anomalía documentada en \S\ref{sec:limitaciones-conocidas}:

\begin{itemize}
\item Recurso en celda propia de MAX (nadie porta, distancia 0): $U = -1 \cdot 0 = 0$
\item MAX porta y está a 4 pasos de su base: $U = +3 - 1\cdot4 = -1$
\item \textbf{Resultado:} Recurso en celda propia ($U=0$) se valora \textbf{mejor} que llevarlo
  a más de 3 pasos de la base ($U < 0$).
\item Además, profundidades pares aplanan valores a 0 (efecto horizonte, caso 10).
\end{itemize}

\textbf{Decisión (2026-10-17):} No cambiar defaults (invalidaría casos 7--10 validados).
Se documenta como limitación; quien necesite otro comportamiento ajusta
\code{recurso\_cargado $\ge$ 8} en el JSON de su escenario.

\section{Explicación para Sustentación: Ejemplo Caso 7}
\label{sec:ejemplo-caso7}

Escenario \code{caso\_07\_decision\_adversarial.json}: mapa 5$\times$7, coste 1, recurso en
(0,1), A1 en (1,0), A2 en (0,2), B1 en (1,2), B2 en (2,4), A mueve.

\begin{enumerate}
\item \textbf{Profundidad 1 (A mueve, evalúa hoja):}
  \begin{itemize}
  \item A2$\to$(0,1) agarra recurso: $U = +3 - 1\cdot1 = 2$ (distancia a base A = 1) + rival term
  \item Valor $\approx$ 3 $\to$ elige agarrar
  \end{itemize}
\item \textbf{Profundidad 2 (B responde):}
  \begin{itemize}
  \item B1 se mueve adyacente a A2 (intercepta): recurso cae, $U$ cae a 0
  \item Valor 0 $\to$ sigue agarrando (empate técnico, tie-break orden)
  \end{itemize}
\item \textbf{Profundidad 3 (A vuelve a mover):}
  \begin{itemize}
  \item A1$\to$(1,1) \textit{protege} (bloquea a B1 antes de que A2 agarre)
  \item B1 no puede interceptar $\to$ A2 agarra y mantiene $\to$ $U = 3$
  \end{itemize}
\item \textbf{Profundidad 4:} A1$\to$(0,0) (entra a base sin recurso) valor 0; B responde...
\end{enumerate}

\textbf{Conclusión:} La utilidad hace que el sistema \textbf{aprenda a proteger} antes de
agarrar, porque la interceptación del rival hace que el valor de agarrar directo caiga a 0.